% Causal importance is not a property of the latent.
% NOTE: ICLR 2027 style file not yet available locally; using article class as a
% placeholder. Swap \documentclass and add \usepackage{iclr2027_conference} when
% the style file lands. Layout will change; no numbers will.
\documentclass[11pt]{article}
\usepackage[margin=1in]{geometry}
\usepackage[T1]{fontenc}
\usepackage{lmodern}          % scalable fonts; microtype expansion needs them
\usepackage{amsmath,amssymb,booktabs,graphicx,microtype,url}
\usepackage[hidelinks]{hyperref}
% GENERATED by src/make_macros.py -- do not edit
\newcommand{\CurveHLNineM}{1.392}
\newcommand{\CurveHLOneTwoM}{1.257}
\newcommand{\CurveHLSixM}{1.098}
\newcommand{\CurveHLThreeM}{1.086}
\newcommand{\CurveSlope}{+0.0270}
\newcommand{\CurveSlopeP}{0.28}
\newcommand{\DepthNLFive}{300}
\newcommand{\DepthNLOneNine}{300}
\newcommand{\DepthNLOneTwo}{300}
\newcommand{\DepthNLTwoFour}{300}
\newcommand{\DepthOrderP}{0.042}
\newcommand{\DepthRatioLFive}{0.384}
\newcommand{\DepthRatioLOneNine}{0.078}
\newcommand{\DepthRatioLOneTwo}{0.360}
\newcommand{\DepthRatioLTwoFour}{0.021}
\newcommand{\DepthSpan}{18}
\newcommand{\GSvsRandHL}{2.269}
\newcommand{\GSvsRandHi}{2.638}
\newcommand{\GSvsRandLo}{1.934}
\newcommand{\GSvsRandN}{240}
\newcommand{\MineN}{240}
\newcommand{\MineVsFrozenHL}{1.192}
\newcommand{\MineVsFrozenHi}{1.379}
\newcommand{\MineVsFrozenLo}{1.026}
\newcommand{\MineVsFrozenP}{0.0213}
\newcommand{\MineVsUntrHL}{2.104}
\newcommand{\MineVsUntrHi}{2.423}
\newcommand{\MineVsUntrLo}{1.826}
\newcommand{\MineVsUntrP}{0.0000}
\newcommand{\NDicts}{9}
\newcommand{\PNormFreq}{14.76}
\newcommand{\PNormRare}{3.60}
\newcommand{\RhoCIHalf}{0.10}
\newcommand{\RhoContrastP}{2.1e-06}
\newcommand{\RhoContrastZ}{4.74}
\newcommand{\RhoDeep}{0.088}
\newcommand{\RhoLFive}{0.358}
\newcommand{\RhoLOneNine}{0.031}
\newcommand{\RhoLOneTwo}{0.339}
\newcommand{\RhoLTwoFour}{0.145}
\newcommand{\RhoShallow}{0.349}
\newcommand{\RhoSlopePerLayer}{-0.0163}
\newcommand{\SDRatioAll}{42.1}
\newcommand{\SDRatioDropOne}{7.2}
\newcommand{\TailMassMax}{73}
\newcommand{\TailMassMin}{36}
\newcommand{\TailMaxMedMax}{464}
\newcommand{\TailMaxMedMin}{33}
\newcommand{\TopFiveMassGS}{65.2}
          % GENERATED by src/make_macros.py -- never edit

\title{Causal Importance Is Not a Property of the Latent}
\author{Valentin No\"el}
\date{}

\begin{document}
\maketitle

\begin{abstract}
% TODO: written last. Anchor: latent identity accounts for \SixArmPctLatent\%
% of variance in a six-arm crossed decomposition; fitting choice and its
% interaction with the latent account for the rest of what isn't residual.
\end{abstract}

\section{Introduction}
A published causal-effect number for a sparse autoencoder (SAE) latent --
Cho et al.\ [2607.20596], Marks et al.\ [2403.19647], Bricken et al.\ [2023],
Templeton et al.\ [2024], Gao et al.\ [2406.04093] -- is reported as if it were
a property of the latent: ablate it, read a KL divergence or logit difference,
publish the number. Table~\ref{tab:audit} audits these five papers, the ones we
could confirm zero-ablate a single latent and read a magnitude rather than a
match-rate accuracy or a multi-latent edit. None report dropping special-token
positions before pooling. None report normalising the readout by the magnitude
of the ablation. None report the near-duplicate structure of the dictionary
they ablate. None report the position within the sequence at which the latent
was measured, or whether that position generalises to others where the same
latent fires.

This is not a checklist paper. The last of those four gaps turns out to be the
largest one, and answering it changes the question. \S\ref{sec:flip} shows that
two of these choices are not merely omissions but can flip the sign of a real
comparison on the same data. \S\ref{sec:property} then asks the question the
audit table implies but no cited paper tests: if the same latent is measured
under a different, equally defensible fitting choice for the dictionary it
came from, does its causal effect rank survive? Six SAE arms sharing a single
initialisation seed, varying only decoder freedom, learning rate, sparsity, and
data order -- exactly the choices a practitioner makes and a reviewer would
accept -- give latent identity \SixArmPctLatent\% of the variance in
log-causal-effect. Arm, latent$\times$arm interaction, and position within
latent give it the other \FigVarianceRest\%. Causal importance, in the sense
the field currently measures it, is predominantly a property of the
measurement, not of the latent.

\section{Related work}
\label{sec:relwork}
\textbf{Sanity Checks for SAEs} (Korznikov, Galichin et al., 2602.14111) show
frozen-random baselines match fully-trained SAEs on reconstruction,
interpretability, sparse probing, and causal editing (0.73 vs 0.72). Their
causal-editing metric is RAVEL's interchange-intervention match-rate accuracy,
not zero-ablation with a magnitude readout, so it is outside this paper's
scope (\S\ref{sec:limitations}); we build our soft-frozen baseline
(\S\ref{sec:setup}) as the construction they describe, but make no claim about
their number. Their synthetic ground-truth result and lazy-training account are
theirs and we do not reclaim them.

\textbf{Are SAEBench Metrics Reliable?} (Chanin et al., 2605.18229) find two
SAEBench metrics, TPP and SCR, fail three independent reliability lenses. Both
metrics zero-ablate a small \emph{set} of latents per concept, not one, so they
sit outside the audit table for a different reason than the RAVEL-based papers
above -- and TPP/SCR are also the only ablation-based metrics in SAEBench, so
their failure is confounded with being the only set-ablation metrics; we do not
attribute it to the artifacts catalogued here. Their mechanism (reseed noise,
low discriminability across training trajectories) is complementary to ours
(directional confounds in a fixed measurement): together the two results say
the ablation-based literature currently cannot distinguish a trained dictionary
from a matched baseline with any of the tools on offer.

\textbf{Are Single-Token SAE Features Causally Necessary?} (Cho et al.,
2607.20596) is the closest published methodology to ours and was investigated
as a reanalysis target. Their headline depth correlation ($\rho=0.81$,
Table 21) survives a sensitivity bound on within-layer latent variance with
room to spare; we did not find grounds to contest it. Independently, our own
implementation of their single-token feature detector, run at their exact
specification (layer 12, Gemma-2-2B), yields $0.41$--$0.44\times$ their
reported prevalence at four times their apparent corpus size -- a stable,
unexplained reproducibility discrepancy we report as a finding about the
detector, not a refutation of their causal claim.

\section{Setup}
\label{sec:setup}
We measure the causal effect of a single SAE latent by ablating its contribution
to the residual stream and reading the change in the model's next-token
distribution. For latent $f$ with decoder direction $d_f$ and activation $a_f$ at
position $t$, we replace $h_t \leftarrow h_t - a_f d_f$ at the SAE's hook point
and record
\begin{equation}
\mathrm{KL}_t = D_{\mathrm{KL}}\!\left(p(\cdot \mid h) \,\|\, p(\cdot \mid h')\right),
\label{eq:kl}
\end{equation}
evaluated at position $t$ and, for the propagation analysis, at $t{+}1$ and
$t{+}3$. Experiments use \texttt{gemma-2-2b} and residual-stream Gemma Scope
dictionaries, a matched trained/soft-frozen pair, and six SAE arms sharing a
single seed (\S\ref{sec:property}); a Gemma-3-1B replication is in progress at
the time of writing.

\paragraph{Estimator.}
Group comparisons use the Hodges--Lehmann estimator, the median of pairwise
ratios, reported with a bootstrap interval. HL is the point estimate belonging to
the Mann--Whitney test; pairing a Mann--Whitney $p$-value with a
ratio-of-medians interval, as we did initially, produces an apparent
contradiction that is an artifact of mixing frameworks.

\paragraph{Convention.}
Ratios of the form $\mathrm{KL}_{t+k}/\mathrm{KL}_t$ are computed
\emph{feature-level}: a per-latent median first, then the ratio of medians across
latents. The observation-level alternative differs by a non-constant factor
$0.97$--$1.35$ and is not used.

\section{One comparison, two conclusions}
\label{sec:flip}
Take the single comparison a practitioner would run first: does this dictionary's
rarest activation-frequency decile carry more or less causal mass than a
frequency-matched decile of an untrained, tied-random dictionary? Under raw KL,
the rare decile is \emph{worse} than random ($\mathrm{HL}=\RawFlipHL$, 95\% CI
$[\RawFlipLo,\RawFlipHi]$, $n=\RawFlipNT$ vs $\RawFlipNR$,
Mann--Whitney $p=\RawFlipP$). Under KL per unit perturbation norm, the identical
latents and the identical control are \emph{better} than random
($\mathrm{HL}=\NormFlipHL$, 95\% CI $[\NormFlipLo,\NormFlipHi]$,
$p=\NormFlipP$). Both intervals exclude one. Nothing about the sample, the model,
or the dictionary changed between the two readings; only whether the
intervention's own size was accounted for did. Perturbation magnitude varies
from $\PNormRare$ to $\PNormFreq$ across activation-frequency deciles
(\S\ref{sec:setup} arms), so raw KL is comparing interventions of systematically
unequal size and mistaking the result for a property of rarity.

\paragraph{Corollary.} Activation frequency predicts causal mass in neither the
trained nor the control arm once magnitude is accounted for. The apparent
frequency dependence in this setting is magnitude, not frequency -- and it
would have shipped as a rarity finding under the more common, uncontrolled
convention.

\section{Causal importance is not a property of the latent}
\label{sec:property}
The comparison in \S\ref{sec:flip} shows that \emph{how} a single fixed dataset
is read can flip a conclusion. This section asks a different question: does the
same latent, measured on dictionaries that differ only in a fitting choice,
keep its rank at all?

\paragraph{Design.} Six TopK SAEs share initialisation seed $0$, so latent $i$
denotes the same initial direction in every arm; each is then trained with one
different, individually defensible choice: decoder free, soft-frozen at
$\tau=0.80$, soft-frozen at $\tau=0.90$, learning rate $10\times$ lower, sparsity
$k=41$ instead of $82$, or a reshuffled corpus order. Every arm sees the same
$12$M tokens. Live latents are intersected across all six arms
($\SixArmSharedN$ shared), one uniform sample is drawn from the intersection,
and every arm is measured on exactly that sample at up to six positions per
latent -- giving a crossed latent$\times$arm design with real replication on
both factors, the direct six-level analogue of a model$\times$wrapper design.

\begin{table}[t]\centering
\begin{tabular}{lcc}
\toprule
component & 2 arms (extrapolated) & 6 arms (direct) \\
\midrule
latent                     & --              & $\SixArmPctLatent\%$ \\
arm                        & $\TwoArmPctArm\%$   & $\SixArmPctArm\%$ \\
latent $\times$ arm        & $\TwoArmPctInter\%$ & $\SixArmPctInter\%$ \\
residual (position)        & --              & $\SixArmPctResid\%$ \\
latent$\times$arm : arm    & $\TwoArmRatio\times$ & $\SixArmRatio\times$ \\
$E\rho^2$ (one-arm design)  & $\TwoArmErhoTwo$    & $\SixArmErhoTwo$ \\
arms needed for $E\rho^2=0.8$ & $\TwoArmArmsNeeded$ (extrapolated) & $\SixArmArmsNeeded$ (direct) \\
\bottomrule
\end{tabular}
\caption{Crossed variance decomposition of log causal effect, latent $\times$
arm, method-of-moments ANOVA on a balanced subsample
($\SixArmI$ latents $\times$ $\SixArmJ$ arms $\times$ $\SixArmN$ positions per
cell). The 2-arm column is the earlier trained-vs-soft-frozen design, with one
degree of freedom on the arm factor; it is superseded by the 6-arm column, not
reported as a separate result.}
\label{tab:decomp}
\end{table}

\paragraph{The headline survives going from one degree of freedom to five.}
The 2-arm design's extrapolated $E\rho^2=\TwoArmErhoTwo$ and its ``$\TwoArmArmsNeeded$
arms needed'' corroborate closely with the 6-arm design's direct estimate
($E\rho^2=\SixArmErhoTwo$, $\SixArmArmsNeeded$ arms needed) -- Table~\ref{tab:decomp}.
This is not an artifact of a single contrast: latent identity accounts for only
$\SixArmPctLatent\%$ of variance, and arm, latent$\times$arm, and position
jointly account for the other $\FigVarianceRest\%$.

\paragraph{What did not survive.} The 2-arm latent$\times$arm-to-arm ratio
($\TwoArmRatio\times$) falls to $\SixArmRatio\times$ with real replication on
the arm factor. A single contrast between two arms is a noisy estimate of the
arm component with one degree of freedom; six defensible fitting choices give
arm a larger, more honest share ($\SixArmPctArm\%$ vs $\TwoArmPctArm\%$). We
report this reversal rather than quietly keeping the larger, earlier number:
the ratio was never the load-bearing claim, and $E\rho^2$ is the number that
held up.

\paragraph{Rank transfers partially, never fully.} Pairwise Spearman
correlation of each latent's causal effect across all $\SixArmPairs$ arm pairs
has median $\SixArmRhoMedian$, range $[\SixArmRhoMin,\SixArmRhoMax]$ -- every
pair positive, none above $0.6$. A latent that is causally important under one
fitting choice is somewhat more likely to be important under another, but a
one-dictionary measurement is closer to a single reviewer's grade than to a
latent's property.

\paragraph{Robustness to the narrowest arm.} The $k{=}41$ arm alone restricts
the shared live-latent intersection to $12{,}080/16{,}384$. Dropping it and
recomputing on the remaining five arms changes every figure by at most $1.2$
points (latent $\SixArmPctLatent\%\to\FiveArmPctLatent\%$, arm
$\SixArmPctArm\%\to\FiveArmPctArm\%$, latent$\times$arm
$\SixArmPctInter\%\to\FiveArmPctInter\%$, residual
$\SixArmPctResid\%\to\FiveArmPctResid\%$, $E\rho^2$
$\SixArmErhoTwo\to\FiveArmErhoTwo$). The decomposition is not driven by the
narrowest arm's dictionary.

\section{Position accounts for most of what is left}
\label{sec:position}
The residual term in Table~\ref{tab:decomp} is position within latent:
$\SixArmPctResid\%$ of variance, larger than latent, arm, and their interaction
combined. Of the five audited papers (Table~\ref{tab:audit}), the ones whose
methodology we could confirm in detail measure each latent at effectively one
position per instance -- a single curated token (Templeton et al.'s case-study
ablations), the top-activating token (the convention we ourselves used before
this decomposition), or an aggregate that is not itself decomposed by position
(Gao et al.'s ablation-sparsity metric). None report what fraction of a
latent's own measured effect is explained by which position was chosen over
another. At $\SixArmPctResid\%$, that choice explains more of the answer than
which SAE the latent came from.

\section{What a defensible measurement needs}
\label{sec:protocol}
Each item below is an artifact established in \S\ref{sec:artifacts}, stated as
the correction it implies. A practitioner reporting ablation-based causal effects
should:
\begin{enumerate}
\item \textbf{Drop special-token positions} before any pooled statistic.
Including position $0$ moves the explained variance of a released Gemma Scope
dictionary from $0.863$ to $-3.5$.
\item \textbf{Report effect per unit perturbation norm}, not raw. \S\ref{sec:flip}
shows this alone reverses the sign of a real comparison.
\item \textbf{Report the per-latent nearest-neighbour cosine distribution}, and
adjust or match on it \emph{in dictionaries that exhibit near-duplicate
structure}. Whether yours does is one line of code.
\item \textbf{Report unembedding alignment} alongside the effect. This is a
sensitivity rather than a confound (\S\ref{sec:artifacts}), so both the
unadjusted and alignment-adjusted estimates belong in the table.
\item \textbf{Report the jackknife curve}, not the mean or the variance ratio.
Removing one latent of $\GSvsRandN$ moves our variance ratio from
$\SDRatioAll\times$ to $\SDRatioDropOne\times$.
\item \textbf{State the sampling scheme.} The size of the alignment sensitivity
depends on it, and no sample of a few hundred latents reaches the extreme tail of
the alignment distribution.
\item \textbf{Report position, and how many positions per latent were measured.}
\S\ref{sec:position}: position accounts for more variance than which dictionary
the latent came from.
\end{enumerate}

\section{Artifacts}
\label{sec:artifacts}

Each subsection gives the mechanism, the evidence, the control, and the effect on
the reported number.

\subsection{Special-token reconstruction}
\textbf{Mechanism.} The BOS residual carries a massive activation the dictionary
does not reconstruct: cosine $0.44$ between input and reconstruction at position
$0$ against $\approx\!0.93$ elsewhere. Its norm dominates the pooled variance.
\textbf{Evidence.} Including position $0$ gives explained variance $-3.5$;
excluding it gives $0.863$ at mean $L_0$ $85.4$, matching the released
specification. \textbf{Control.} Drop special-token positions; assert the
expected explained variance and abort otherwise. Without that assertion every
downstream number is computed on a broken hook point.

\subsection{Perturbation magnitude}
\textbf{Mechanism.} Ablation perturbs the residual stream by $a_f d_f$, and
activation magnitude covaries with activation frequency, so raw effect sizes
compare interventions of unequal size. \textbf{Evidence.} \S\ref{sec:flip}.
Median perturbation norm runs $\PNormRare$ in the rarest frequency decile to
$\PNormFreq$ in the most frequent. \textbf{Control.} Report KL per unit
perturbation norm throughout.

\subsection{Near-duplicate structure}
\textbf{Mechanism.} Feature splitting produces decoder rows at high mutual
cosine; ablating one member of such a pair leaves the concept partly represented.
\textbf{Evidence.} Per-latent maximum cosine predicts causal mass in the released
Gemma Scope dictionary ($\rho=+0.220$, $p=0.0006$) and in neither of our trained
arms ($+0.078$, $+0.060$), which have visibly less splitting (medians $0.207$ and
$0.122$ against $0.272$). \textbf{Control.} One-sided adjustment where the
relationship is asymmetric. Symmetric residualisation is wrong here: fitting the
control arm produces a slope of $+2.506$ against a predictor with $p=0.82$, i.e.\
residualising against noise. Matching is infeasible against a random dictionary,
which structurally cannot contain near-duplicates.

\paragraph{This artifact is conditional.} Near-duplicate structure confounds
causal-effect measurement \emph{in dictionaries that have it}. That is more
useful than a universal claim, because it comes with its own diagnostic.

\subsection{Direct-path alignment: a sensitivity, not a confound}
\label{sec:align}
\textbf{Mechanism.} A trained decoder direction is aligned with directions the
model writes to the residual stream, and those feed the unembedding. Next-token
KL therefore rewards latents that short-circuit the network. We measure alignment
as $\sqrt{d^\top C d}$ with $C$ the embedding covariance --- the spread of the
direct logit contribution, since a uniform logit shift leaves the softmax
unchanged.

\textbf{Evidence.} Alignment predicts causal mass at $\rho=+0.375$ in Gemma Scope
and $+0.264$ in a dictionary we trained ourselves with a different architecture
and roughly $40\times$ less data, with the control arm null in both cases
($-0.063$, $+0.004$). The single most causally powerful latent we found writes the
letter \emph{k} to the logits and carries $12\times$ the causal mass of the
runner-up.

\textbf{Why this is a sensitivity and not a confound.} The trivially aligned
latents are approximately the top $20$ of $16{,}384$, the $99.88$th percentile. No
sample a paper would realistically draw reaches them: the top twenty by alignment
within a $240$-latent sample sit near the $92$nd dictionary percentile and are
$90\%$ word-like, under both stratified and uniform sampling. So at realistic
sample sizes high-alignment latents are semantic, and adjusting alignment away
removes signal rather than triviality. The correct statement is that
ablation-based measurements are sensitive to unembedding alignment at
$\rho\approx0.3$, and that within realistic samples this alignment reflects
semantic content.

\textbf{An unexplained depth dependence.} Because $C$ is layer-independent by
construction, the measure estimates the logit effect a direction \emph{would} have
if it reached the output unchanged, so it should predict best where fewest blocks
intervene --- deepest. It does the opposite. Alignment predicts causal mass with
$\rho=\RhoLFive$ at layer 5 and $\RhoLOneTwo$ at layer 12, against $\RhoLOneNine$
at layer 19 and $\RhoLTwoFour$ at layer 24 ($n=\DepthNLFive$ per depth).

At $n=300$ a Spearman correlation of $0.3$ carries a $95\%$ interval of about
$\pm\RhoCIHalf$, so the two shallow estimates are indistinguishable from each
other and so, marginally, are the two deep ones. We therefore test the trend
rather than describe the sequence: pooling by Fisher $z$ gives
$\rho=\RhoShallow$ for the shallow pair and $\rho=\RhoDeep$ for the deep pair, a
contrast of $z=\RhoContrastZ$ ($p=\RhoContrastP$), and a weighted regression of
$z$ on layer index has slope $\RhoSlopePerLayer$ per layer where the measure's
construction predicts a positive slope. The pattern is a decrease with depth, not
a mid-network dip with recovery.

The compressed-variance explanation is falsified: alignment spread is equal across
depths (CV $0.1175$ against $0.1091$) and causal-mass spread is \emph{greater} at
layer 19. We report the dependence as unexplained; it is a violation of the
measure's own construction rather than an observation about the model.

\subsection{Position in sequence}
\textbf{Mechanism.} Effects are measured at each latent's top-activating
positions, and later positions carry more context and a sharper next-token
distribution. \textbf{Evidence.} Pooled $\rho=-0.087$ ($p<10^{-4}$) across
$2{,}700$ latents, reaching $-0.207$ in one dictionary and null in five. The sign
is \emph{negative}: earlier positions carry slightly more causal mass, opposite to
the mechanism we proposed. On the endpoint sample the effect is not significant
($\rho=-0.084$, $p=0.196$). This correlational null does not contradict
\S\ref{sec:position}'s variance share: a latent's effect can vary substantially
across its own top positions ($\SixArmPctResid\%$ of variance) without that
variation trending consistently early-to-late across latents ($\rho\approx0$).
\textbf{Control.} Record and report position, and how many positions per latent
were sampled; a single position per latent cannot distinguish a stable causal
effect from a coincidence of which token was chosen.

\section{Results}
\label{sec:results}

\subsection{Endpoints}
Table~\ref{tab:endpoints} gives the trained-versus-control comparison for three
controls. The second and third rows share a dictionary, a feature list and an
evaluation, so their ratio is a within-experiment decomposition.

\begin{table}[t]\centering
\begin{tabular}{lccc}
\toprule
comparison & HL & 95\% CI & $n$ \\
\midrule
Gemma Scope vs untrained tied & $\GSvsRandHL$ & $[\GSvsRandLo, \GSvsRandHi]$ & $\GSvsRandN$ \\
ours vs untrained tied        & $\MineVsUntrHL$ & $[\MineVsUntrLo, \MineVsUntrHi]$ & $\MineN$ \\
ours vs soft-frozen           & $\MineVsFrozenHL$ & $[\MineVsFrozenLo, \MineVsFrozenHi]$ & $\MineN$ \\
\bottomrule
\end{tabular}
\caption{Causal effect per unit perturbation norm, Hodges--Lehmann ratios on
uniform samples of live latents. The soft-frozen control has a trained encoder
and a decoder constrained to cosine $0.8$ of its random initialisation.}
\label{tab:endpoints}
\end{table}

A trained dictionary beats an untrained one by $\MineVsUntrHL\times$ and a
soft-frozen baseline by $\MineVsFrozenHL\times$ ($p=\MineVsFrozenP$). Most of the
measured advantage is therefore attributable to the encoder; decoder freedom
contributes the smaller factor. This holds despite the two arms differing sharply
in decoder geometry --- median cosine to initialisation $0.231$ for the free
decoder against $0.800$ for the constrained one --- and despite a reconstruction
gap that grows monotonically over training ($+0.0064$ explained variance per
million tokens, $p<10^{-4}$).

\subsection{The tail is not a population phenomenon}
Across $\NDicts$ Gemma Scope dictionaries spanning two depths, three widths and
sparsity from $L_0$ $22$ to $445$, the ratio of maximum to median causal mass
within a fixed sample of $300$ latents runs from $\TailMaxMedMin\times$ to
$\TailMaxMedMax\times$, and the top $5\%$ of latents carry between
$\TailMassMin\%$ and $\TailMassMax\%$ of total causal mass. Every dictionary shows
it.

But within any one dictionary the tail is a handful of latents. Removing the
single largest takes the variance ratio from $\SDRatioAll\times$ to
$\SDRatioDropOne\times$; the median ratio is unmoved. A bootstrap resamples those
latents in and out while still treating them as population draws, so the jackknife
curve is the honest instrument. At $n\approx200$ the control arm's $99$th
percentile has a bootstrap interval spanning $11.3\times$, and the derived
``fraction above the control's $p99$'' runs $[4.2\%, 39.6\%]$: published aggregate
comparisons on a few hundred latents cannot distinguish anything.

\subsection{Effect decays with remaining depth, but not proportionally}
Measuring \eqref{eq:kl} at $t$ and $t{+}3$ across four depths at matched width and
approximately matched sparsity gives $\mathrm{KL}_{t+3}/\mathrm{KL}_t$ of
$\DepthRatioLFive$ (layer 5, $20$ blocks remaining), $\DepthRatioLOneTwo$ (layer 12,
$13$), $\DepthRatioLOneNine$ (layer 19, $6$) and $\DepthRatioLTwoFour$ (layer 24, $1$),
with $n=\DepthNLFive$ latents per depth. The effect itself is large and not in
doubt: a span of $\DepthSpan\times$ from shallowest to deepest.

The \emph{attribution} to propagation limits is much weaker, and we state it as
such. We pre-registered three predictions of decreasing strength. A power law in
remaining depth, fitted from layers 12 and 19 and tested out of sample with a
factor-of-$1.5$ tolerance, predicts $0.760$ at layer 5 and $0.003$ at layer 24
against observed $\DepthRatioLFive$ and $\DepthRatioLTwoFour$ --- misses of
$0.51\times$ and $8.4\times$, in opposite directions, so the curve is flatter than
proportional at both ends. Two threshold predictions, $>\!0.5$ at layer 5 and
$<\!0.02$ at layer 24, also both failed. Only bare monotone ordering survived, and
on four points that is worth $p\approx\DepthOrderP$ under a random-permutation
null.

So: $\mathrm{KL}_{t+3}/\mathrm{KL}_t$ declines with depth by a factor of
$\DepthSpan$; the propagation account predicts this ordinally at
$p\approx\DepthOrderP$ and is falsified as a proportionality in both directions.
The decline is solid, the mechanism is not established.

\paragraph{Consequence for \S\ref{sec:align}.} Because this gradient is a
propagation effect, it is not evidence that deeper layers are more direct-path
dominated, and it must not be read as support for the alignment sensitivity.

\subsection{Reconstruction improves without a matching causal gain}
Evaluating the trained and soft-frozen arms at four training budgets gives HL
ratios of $\CurveHLThreeM$, $\CurveHLSixM$, $\CurveHLNineM$ and $\CurveHLOneTwoM$ at $3$, $6$,
$9$ and $12$ million tokens (slope $\CurveSlope$ per million, $p=\CurveSlopeP$),
while the explained-variance gap over the same range grows from $0.000$ to
$0.077$ at $p<10^{-4}$. The reconstruction gap grows significantly and the causal
gap does not. With four budgets, $120$ latents each, and non-monotone point
estimates, this is consistent with a dissociation but does not establish one.

\section{Errors we made}
\label{sec:errors}
We log these because each would have shipped as a real result if not caught, and
because the reversals are themselves informative about where this kind of
measurement goes wrong.

\textbf{A sign error made a control column constant.} Zeroing a self-similarity
diagonal after taking an absolute value, rather than before, left every entry
in one column equal to $2.0$ instead of excluded. The near-duplicate analysis
in \S\ref{sec:artifacts} was rebuilt after this was caught.

\textbf{Residualising against noise fabricated a confound.} Fitting a
covariate-adjustment regression on the random control arm -- whose predictor
has no true relationship to the outcome ($p=0.82$) -- produced a large,
significant-looking slope ($+2.506$) that appeared to explain away a real
effect in the trained arm. Adjustment must be asymmetric: fit only where the
covariate relationship is real.

\textbf{Pairing a Mann--Whitney $p$-value with the wrong interval created an
apparent contradiction.} A ratio-of-medians bootstrap CI containing $1$
alongside a significant Mann--Whitney test looks contradictory because the two
statistics belong to different estimators. Hodges--Lehmann is the point
estimate the Mann--Whitney test actually tests; \S\ref{sec:setup}'s convention
was adopted after this was diagnosed, not before.

\textbf{A falsified prediction revealed the real mechanism.} We predicted
uniform (rather than frequency-stratified) sampling would show top-alignment
latents to be geometrically trivial. It did not: wordlike fraction was
$0.90$ under both schemes. The real reason high-alignment latents look
semantic is sample size, not sampling scheme -- the top alignment percentile is
unreachable at realistic $n$ regardless of how you sample. The falsification is
what produced \S\ref{sec:align}'s actual argument.

\textbf{An overclaimed headline was walked back to what survived.} An early
draft of \S\ref{sec:results}'s depth result stated ``monotonicity confirmed'' as
the primary finding. Of three pre-registered predictions, only the weakest
(bare ordering, $p\approx\DepthOrderP$ on four points) survived; the power law
and both thresholds were falsified. The corrected framing --- decline is solid,
mechanism is not established --- is what appears above.

\section{Limitations and conclusion}
\label{sec:limitations}
\textbf{One model, for now.} All measurements here use \texttt{gemma-2-2b}; a
Gemma-3-1B replication of the six-arm decomposition is in progress. Whether the
depth dependence in \S\ref{sec:align} is about relative depth or about this
model is separately untested; Llama Scope dictionaries would answer that.
\textbf{Training budget.} Our trained arms see $12$M tokens against Gemma
Scope's billions, and the reconstruction gap is still widening at that point,
so any claim about what decoder training does is scoped to this budget.
\textbf{One corpus} for the trained arms (WikiText-103), and one evaluation
corpus. \textbf{Ablation only.} These artifacts apply to metrics that zero or
subtract a latent's contribution and read a magnitude. They do not transfer
directly to interchange-based scores such as RAVEL, which are match-rate
accuracies; we make no claim about that result (\S\ref{sec:relwork}). Ablation, interchange and
clamping all perturb along the same $d_f$, so ``use interchange instead'' is
not a remedy --- the axis that matters is magnitude versus rate.
\textbf{Attribution not established.} The two SAEBench metrics that fail
reliability auditing are the only ablation-based metrics in that suite, but also
the only set-ablation metrics; the properties are confounded and we do not claim
their failure follows from these artifacts. \textbf{Unexplained.} The depth
dependence of the alignment sensitivity, and the sign of the position-trend
effect in \S\ref{sec:artifacts} (distinct from the position-\emph{variance}
result in \S\ref{sec:position}, which is about spread within a latent rather
than a trend across latents).

Five audited papers report a single-latent causal effect with a magnitude
readout; none report the position, duplication, or magnitude controls this
paper establishes as necessary, and a sixth family of closely related papers
(RAVEL-based, set-ablation) uses a different readout entirely for reasons we
make explicit rather than dispute. Two of the controllable choices flip the
sign of a real comparison (\S\ref{sec:flip}). A crossed six-arm design --
the first, to our knowledge, to give the arm factor real replication in this
setting -- shows that a latent's own identity accounts for
$\SixArmPctLatent\%$ of the variance in its measured causal effect, and that
which position of that latent's own activations was measured accounts for
more than that. Causal importance, as the field currently measures it, is a
property of the measurement first and the latent second.

\begin{table}[t]\centering
\footnotesize
\begin{tabular}{p{2.1cm}p{3.6cm}p{1.5cm}p{1.5cm}p{1.9cm}p{1.9cm}}
\toprule
paper & readout & drops special tok. & norm.\ by magnitude & reports duplication & accounts for alignment \\
\midrule
Bricken et al.\ 2023 & per-feature ablation; output logit change & --- & --- & broad, not tied to effect size & --- \\
Marks et al.\ 2024 & single-feature zero-ablation; KL divergence & --- & --- & --- & --- \\
Templeton et al.\ 2024 & per-feature zero-clamp vs.\ linear attribution & --- & --- & broad, not tied to effect size & --- \\
Gao et al.\ 2024 & per-latent ablation; sparsity of effect vector (not magnitude) & --- & different axis & --- & --- \\
Cho et al.\ 2025 & single-latent ablation; KL/logit necessity score & this paper's Artifact 1 & --- & --- & unclear vs.\ $\sqrt{d^\top C d}$ \\
\bottomrule
\end{tabular}
\caption{Audit of the five papers we could confirm use single-latent
zero-ablation with a magnitude readout. ``---'' means not stated in the
sections we reviewed, not confirmed absent. Excluded as a different readout
family: Makelov et al.\ 2024 (multi-feature edits), SAEBench TPP/SCR (set
ablation), Korznikov/Galichin et al.\ 2026 (RAVEL match-rate), Rajamanoharan
et al.\ 2024 (whole-dictionary substitution) -- see \S\ref{sec:relwork}.}
\label{tab:audit}
\end{table}

\bibliographystyle{plain}
% \bibliography{refs}

\end{document}
